% Included by main.tex; the preamble and references are shared.
\section{A variational principle for tangent vector fields}\label{sec:vector-fields}

For $X\in H^1(\Omega;\mathbb R^2)$, \emph{tangent trace} means
$\operatorname{Tr}X\cdot\nu=0$ almost everywhere on $\partial\Omega$;
see \cite[Theorem 1.1, p.~2]{Rohleder} for this boundary condition.
For any $X\in H^1(\Omega;\mathbb R^2)$ define
\[
 Q(X)=\int_\Omega\bigl(|\operatorname{div}X|^2
                         +|\operatorname{curl}X|^2\bigr),
 \qquad \operatorname{curl}X=\partial_1X_2-\partial_2X_1.
\]
The following is the divergence--curl form of Rohleder's vector
variational principle \cite[Theorem 1.1]{Rohleder}; see also \cite{Roh24}.
We give its derivation,
including the equality case.

\begin{proposition}\label[proposition]{prop:vector-variational}
For every $X\in H^1(\Omega;\mathbb R^2)$ with tangent trace,
\begin{equation}\label{eq:vector-inequality}
 Q(X)\ge\lambda_N\|X\|_2^2.
\end{equation}
Equality holds if and only if $X=\nabla v$ for some $v\in V$.
\end{proposition}
\begin{proof}
We use the Sobolev integration-by-parts identity
\begin{equation}\label{eq:sobolev-ibp}
 \int_\Omega Y\cdot\nabla u
 =-\int_\Omega u\,\operatorname{div}Y
   +\int_{\partial\Omega}\operatorname{Tr}u\,
        (\operatorname{Tr}Y\cdot\nu)\,d\sigma,
\end{equation}
valid for $u\in H^1(\Omega;\mathbb R)$ and
$Y\in H^1(\Omega;\mathbb R^2)$. It follows from the classical divergence
theorem by smooth approximation in $H^1$ and continuity of the trace
into $L^2(\partial\Omega)$; see \cite[Theorem 1.46, Theorem 3.44,
and Section 3.11]{Hunter}. The boundary product is integrable by
Cauchy--Schwarz.

First solve
\begin{equation}\label{eq:potentials}
 \begin{aligned}
 \Delta v&=\operatorname{div}X,\qquad
 \partial_\nu v=0,\qquad \int_\Omega v=0,\\
 \Delta j&=\operatorname{curl}X,\qquad \operatorname{Tr}j=0.
 \end{aligned}
\end{equation}
For the Neumann problem, the required compatibility condition is
\[
 \int_\Omega\operatorname{div}X
 =\int_{\partial\Omega}\operatorname{Tr}X\cdot\nu\,d\sigma=0.
\]
Here we applied \cref{eq:sobolev-ibp} with $u=1$ and $Y=X$.
On the closed mean-zero subspace
\[
 H=\left\{w\in H^1(\Omega;\mathbb R):\int_\Omega w=0\right\},
\]
the form $a(v,w)=\int_\Omega\nabla v\cdot\nabla w$ is continuous and
coercive by the mean-zero Poincar\'e inequality. Since
$\operatorname{div}X\in L^2(\Omega)$, Lax--Milgram
\cite[Theorem 4.20, p.~103]{Hunter} gives a unique $v\in H$ such that
\[
 \int_\Omega\nabla v\cdot\nabla w
 =-\int_\Omega(\operatorname{div}X)w\qquad(w\in H).
\]
Compatibility extends this identity to all $w\in H^1(\Omega;\mathbb R)$
by subtracting their means. It is the weak formulation of the Neumann
problem in \cref{eq:potentials}.
Similarly, the Dirichlet Poincar\'e inequality makes $a$ coercive on
$H_0^1(\Omega;\mathbb R)$, so Lax--Milgram gives a unique
$j\in H_0^1(\Omega;\mathbb R)$ satisfying
\[
 \int_\Omega\nabla j\cdot\nabla w
 =-\int_\Omega(\operatorname{curl}X)w
 \qquad(w\in H_0^1(\Omega;\mathbb R)).
\]
No compatibility condition is needed for this Dirichlet problem.
Boundary regularity gives $v,j\in H^2(\Omega)$; see
\cite[Chapter III, Theorems 6.1 and 6.4]{Showalter}.

Put $\nabla^\perp j=(-\partial_2j,\partial_1j)$.
Tangential differentiation commutes with taking the trace for $H^2$
functions on a smooth domain; this follows by smooth approximation and
continuity of the trace maps. Thus the zero trace of $j$ implies that
$\nabla^\perp j$ has tangent trace. Consequently
\[
 Z=X-\nabla v-\nabla^\perp j
\]
belongs to $H^1(\Omega;\mathbb R^2)$ and has tangent trace.
The equations in \cref{eq:potentials} give
$\operatorname{div}Z=\operatorname{curl}Z=0$. Its components are
harmonic in distributions and hence smooth in the interior by Weyl's
lemma. Simple connectivity therefore gives a smooth real interior
potential $r$ with $Z=\nabla r$.

For $M>0$, define $r_M=\max\{-M,\min\{r,M\}\}$, and write
$s^+=\max\{s,0\}$.
The weak derivative formulas for positive parts
\cite[Proposition 3.22 and the subsequent discussion, pp.~57--58]{Hunter},
applied locally to
$r_M=r-(r-M)^++(-r-M)^+$, give
\[
 r_M\in H^1_{\mathrm{loc}}(\Omega),\qquad
 \nabla r_M=\mathbf1_{\{|r|<M\}}Z\quad\text{almost everywhere}.
\]
Since $|r_M|\le M$, $\Omega$ is bounded, and $Z\in L^2(\Omega)$,
it follows that $r_M\in H^1(\Omega)$.

Fix a ball $B$ with $\overline B\subset\Omega$, and write $u_B$ and
$u_\Omega$ for the averages over $B$ and $\Omega$. The usual mean-zero
Poincar\'e inequality gives, for $u\in H^1(\Omega)$,
\[
 \begin{aligned}
 |u_B-u_\Omega|
 &\le |B|^{-1/2}\|u-u_\Omega\|_2,\\
 \|u-u_B\|_2
 &\le\left(1+\sqrt{\frac{|\Omega|}{|B|}}\right)
       \|u-u_\Omega\|_2\\
 &\le C_{\Omega,B}\|\nabla u\|_2.
 \end{aligned}
\]
Applying this to $r_M$ and using $|r_M|\le|r|$ on $B$ yields
\[
 \|r_M\|_2
 \le C_{\Omega,B}\|Z\|_2
       +|\Omega|^{1/2}|(r_M)_B|
 \le C_{\Omega,B}\|Z\|_2
       +|\Omega|^{1/2}\sup_B|r|.
\]
The bound is independent of $M$. Since $r_M\to r$ pointwise, Fatou's
lemma gives $r\in L^2(\Omega)$. Moreover, for every
$\phi\in C_c^\infty(\Omega)$,
\[
 \int_\Omega r\,\partial_i\phi=-\int_\Omega Z_i\phi
 \qquad(i=1,2),
\]
by the interior identity $\nabla r=Z$. Thus the global weak gradient of
$r$ is $Z$, and $r\in H^1(\Omega)$. This establishes its global
integrability using only the interior potential and $Z\in L^2$.
Applying \cref{eq:sobolev-ibp} with $u=r$ and $Y=Z$ now gives
\[
 \int_\Omega|\nabla r|^2
 =\int_\Omega Z\cdot\nabla r
 =-\int_\Omega r\,\operatorname{div}Z
   +\int_{\partial\Omega}\operatorname{Tr}r\,
        (\operatorname{Tr}Z\cdot\nu)\,d\sigma=0.
\]
Hence $Z=0$, and
\begin{equation}\label{eq:decomposition}
 X=\nabla v+\nabla^\perp j.
\end{equation}
The summands are orthogonal in $L^2$, since
\[
 \int_\Omega\nabla v\cdot\nabla^\perp j
 =-\int_\Omega v\,\operatorname{div}(\nabla^\perp j)
   +\int_{\partial\Omega}\operatorname{Tr}v\,
       (\operatorname{Tr}\nabla^\perp j\cdot\nu)\,d\sigma=0.
\]
Consequently
\[
 \|X\|_2^2=\|\nabla v\|_2^2+\|\nabla j\|_2^2,
 \qquad Q(X)=\|\Delta v\|_2^2+\|\Delta j\|_2^2.
\]

Integration by parts, Cauchy--Schwarz, and the two Poincar\'e inequalities
give
\[
 \begin{aligned}
 \|\nabla v\|_2^2
 &=-\int_\Omega v\,\Delta v
 \le\|v\|_2\|\Delta v\|_2\\
 &\le\lambda_N^{-1/2}\|\nabla v\|_2\|\Delta v\|_2,\\
 \|\nabla j\|_2^2
 &=-\int_\Omega j\,\Delta j
 \le\|j\|_2\|\Delta j\|_2\\
 &\le\lambda_D^{-1/2}\|\nabla j\|_2\|\Delta j\|_2.
 \end{aligned}
\]
The boundary terms vanish because $\partial_\nu v=0$ and
$\operatorname{Tr}j=0$, respectively. The last inequalities use the
zero mean of $v$ and the zero trace of $j$. Dividing by the gradient
norms when nonzero, with the zero cases immediate, gives
\[
 \|\Delta v\|_2^2\ge\lambda_N\|\nabla v\|_2^2,
 \qquad
 \|\Delta j\|_2^2\ge\lambda_D\|\nabla j\|_2^2.
\]
It follows that
\begin{equation}\label{eq:vector-refined}
 \begin{aligned}
 Q(X)&\ge\lambda_N\|\nabla v\|_2^2
              +\lambda_D\|\nabla j\|_2^2\\
 &=\lambda_N\|X\|_2^2
              +(\lambda_D-\lambda_N)\|\nabla j\|_2^2.
 \end{aligned}
\end{equation}
This proves \cref{eq:vector-inequality}.

Suppose equality holds in \cref{eq:vector-inequality}.
Since $\lambda_D>\lambda_N$, \cref{eq:vector-refined} forces
$\nabla j=0$, and the zero trace gives $j=0$. Thus $X=\nabla v$ and
$\|\Delta v\|_2^2=\lambda_N\|\nabla v\|_2^2$.
If $\nabla v\ne0$, the preceding chain becomes
\[
 \begin{aligned}
 \|\nabla v\|_2^2
 &\le\|v\|_2\|\Delta v\|_2\\
 &\le\lambda_N^{-1/2}\|\nabla v\|_2\|\Delta v\|_2
 =\|\nabla v\|_2^2.
 \end{aligned}
\]
Both inequalities are equalities, so
$\|\nabla v\|_2^2=\lambda_N\|v\|_2^2$.
The equality characterization of the mean-zero Neumann variational
inequality (Fact~2.2 of the spectral and planar preliminaries) gives
\cref{eq:V-weak}, hence $v\in V$.
If $\nabla v=0$, its zero mean gives $v=0\in V$.

Conversely, if $v\in V$, then $X=\nabla v$ belongs to
$H^1(\Omega;\mathbb R^2)$ and has tangent trace because
$\partial_\nu v=0$. Also $\operatorname{curl}X=0$ and
$\operatorname{div}X=\Delta v=-\lambda_Nv$, so
\[
 Q(X)=\lambda_N^2\|v\|_2^2
      =\lambda_N\|\nabla v\|_2^2
      =\lambda_N\|X\|_2^2,
\]
where the middle equality follows by testing \cref{eq:V-weak} with $v$.
\end{proof}

For a differential-form interpretation of the divergence--curl energy,
see \cite[Section 5.1]{FGM26}.

We next record the boundary identity in fixed Cartesian coordinates.
Let $\tau=(-\nu_2,\nu_1)$ be the counterclockwise unit tangent, and set
\[
 \kappa=\det(\tau,\partial_\tau\tau).
\]
Thus $\kappa$ is positive on a convex circle. Here
$|\nabla X|^2=\sum_{i,j=1}^2(\partial_iX_j)^2$.

\begin{lemma}\label[lemma]{lem:curvature-identity}
For every $X\in H^1(\Omega;\mathbb R^2)$ with tangent trace,
\begin{equation}\label{eq:curvature-identity}
 Q(X)=\int_\Omega|\nabla X|^2
       +\int_{\partial\Omega}\kappa|\operatorname{Tr}X|^2\,d\sigma.
\end{equation}
\end{lemma}
\begin{proof}
First suppose $X\in C^\infty(\overline\Omega;\mathbb R^2)$, without
imposing tangency. Expanding the bulk integrands gives
\[
 \begin{aligned}
 &|\operatorname{div}X|^2+|\operatorname{curl}X|^2-|\nabla X|^2\\
 &\quad=(\partial_1X_1+\partial_2X_2)^2
        +(\partial_1X_2-\partial_2X_1)^2\\
 &\qquad\quad-\sum_{i,j=1}^2(\partial_iX_j)^2\\
 &\quad=2\bigl(\partial_1X_1\,\partial_2X_2
              -\partial_1X_2\,\partial_2X_1\bigr)
        =2\det\nabla X.
 \end{aligned}
\]
Put
\[
 F=\bigl(X_1\partial_2X_2-X_2\partial_2X_1,\,
         X_2\partial_1X_1-X_1\partial_1X_2\bigr).
\]
The mixed second derivatives cancel, so
$\operatorname{div}F=2\det\nabla X$. On the boundary,
\[
 F\cdot\nu=X_1\partial_\tau X_2-X_2\partial_\tau X_1.
\]
The divergence theorem therefore gives
\begin{equation}\label{eq:smooth-boundary-pairing}
 Q(X)-\int_\Omega|\nabla X|^2
 =\int_{\partial\Omega}
   \bigl(X_1\partial_\tau X_2-X_2\partial_\tau X_1\bigr)\,d\sigma.
\end{equation}

For $h=(h_1,h_2)\in H^{1/2}(\partial\Omega;\mathbb R^2)$, define
\[
 B(h)=\langle\partial_\tau h_2,h_1\rangle
       -\langle\partial_\tau h_1,h_2\rangle,
\]
where the brackets denote the real $H^{-1/2}$--$H^{1/2}$ duality
pairing, agreeing with the boundary integral for smooth functions.
Tangential differentiation is bounded from $H^{1/2}(\partial\Omega)$
to $H^{-1/2}(\partial\Omega)$, so $B$ is continuous on $H^{1/2}$.
Choose $X_n\in C^\infty(\overline\Omega;\mathbb R^2)$ converging to
$X\in H^1(\Omega;\mathbb R^2)$ in $H^1$. Their traces converge to
$h=\operatorname{Tr}X$ in $H^{1/2}$, hence
$B(\operatorname{Tr}X_n)\to B(h)$. The bulk integrals also converge,
since the first derivatives converge in $L^2$. Passing to the limit in
\cref{eq:smooth-boundary-pairing} gives
\begin{equation}\label{eq:sobolev-boundary-pairing}
 Q(X)-\int_\Omega|\nabla X|^2=B(h).
\end{equation}
The approximating fields need not have tangent trace.

Now suppose $h$ is tangent, and put $a=h\cdot\tau$, so $h=a\tau$.
Multiplication by the smooth components of $\tau$ preserves $H^{1/2}$,
and hence $a\in H^{1/2}(\partial\Omega;\mathbb R)$.
First take $a$ smooth. The product rule gives
\[
 \begin{aligned}
 B(a\tau)
 &=\int_{\partial\Omega}
   \bigl[a\tau_1\partial_\tau(a\tau_2)
         -a\tau_2\partial_\tau(a\tau_1)\bigr]\,d\sigma\\
 &=\int_{\partial\Omega}
       a^2\det(\tau,\partial_\tau\tau)\,d\sigma
 =\int_{\partial\Omega}\kappa a^2\,d\sigma,
 \end{aligned}
\]
because the terms containing $a\,\partial_\tau a$ cancel.
For general $a\in H^{1/2}$, choose smooth $a_n\to a$ in $H^{1/2}$.
Then $a_n\tau\to a\tau$ in $H^{1/2}$, so continuity of $B$ gives
$B(a_n\tau)\to B(a\tau)$. Also $a_n\to a$ in $L^2$, and
\[
 \|a_n^2-a^2\|_{L^1(\partial\Omega)}
 \le\|a_n-a\|_{L^2(\partial\Omega)}
      \bigl(\|a_n\|_{L^2(\partial\Omega)}
            +\|a\|_{L^2(\partial\Omega)}\bigr)\to0.
\]
Since $\kappa$ is bounded, we may pass to the limit to obtain
$B(a\tau)=\int_{\partial\Omega}\kappa a^2\,d\sigma$.
Finally $|h|^2=a^2$, and \cref{eq:sobolev-boundary-pairing} proves
\cref{eq:curvature-identity}.
\end{proof}
